% ============================================================
%  Поглавље 4: Класични приступ — Z3 решавач
% ============================================================
\section{Класични приступ: Z3 решавач и SMT}
\label{sec:z3}

% --- Увод у SAT/SMT решаваче ---
% Кратко представите SAT решаваче (DPLL, CDCL), а затим SMT
% (Satisfiability Modulo Theories) као проширење са додатним теоријама
% (аритметика, низови, битвектори...).

\subsection{SAT и SMT решавачи}

% TODO: Опишите DPLL/CDCL архитектуру у 1--2 параграфа.
% Поменути Nelson-Oppen оквир за комбиновање теорија у SMT.

\subsection{Z3: Архитектура и могућности}

% --- Опишите Z3 ---
% Z3 је Microsoft Research SMT решавач (de Moura & Bjørner, 2008).
% Подржава: линеарну аритметику, низове, битвекторе, квантификаторе...
% Сучеља: Python API (z3-solver), C, .NET.

% TODO: Укључити кратак блок Python кода са z3-solver:
% deklaracija varijabli, dodavanje ograničenja, poziv solve().

\subsection{Кодирање MaxCut-а у Z3}

% --- Кодирање ---
% За граф G=(V,E), уведите булове променљиве x_v ∈ {0,1} за сваки чвор.
% Ивица {u,v} прелази рез ако x_u ≠ x_v, тј. x_u XOR x_v = 1.
% Циљ: максимизовати Σ_{(u,v)∈E} (x_u XOR x_v).
%
% SMT кодирање: Int или Bool за x_v, Optimize() за максимизацију.

\begin{example}[MaxCut на $K_4$ у Z3 (Python)]
\begin{verbatim}
from z3 import *
x = [Int(f'x{i}') for i in range(4)]
opt = Optimize()
for v in x:
    opt.add(Or(v == 0, v == 1))
edges = [(0,1),(0,2),(0,3),(1,2),(1,3),(2,3)]
cut = Sum([If(x[u] != x[v], 1, 0) for u,v in edges])
opt.maximize(cut)
print(opt.check(), opt.model())
\end{verbatim}
\end{example}

% TODO: Показати излаз за K4 (оптимум = 4) и прокоментарисати.

\subsection{Ограничења SMT приступа}

% --- Скалабилност ---
% Употреба Z3 за MaxCut је прецизна (даје оптимум), али NP-тешкоћа
% проблема значи да се време извршавања у најгорем случају
% експоненцијално повећава са бројем чворова.
% Приложити табелу: n=10,20,50 чворова vs. CPU време у Z3.

% TODO: Мерења перформанси Z3 за случајне графове растуће величине.
% Ово мотивише прелаз на квантни приступ.

\begin{remark}
Z3 је погодан алат за верификацију малих инстанци и за генерисање
референтних решења, али не може заменити апроксимационе или квантне
алгоритме за велике улазе.
\end{remark}
